% ECONOMICS_PROBLEM: {"id": "D86-21", "title": "Strongly polynomial computation of an optimal finite contract", "jel_code": "D86", "source_id": "OP-0596"}
% FIRST_POSTED: 2026-10-09
% ATTEMPT_STATUS: unresolved
% ENTRY_KIND: research_attempt
\documentclass[11pt]{article}
\usepackage[T1]{fontenc}
\usepackage{amsmath,amssymb,amsthm,geometry,hyperref}
\geometry{margin=1in}
\newtheorem{theorem}{Theorem}
\newtheorem{lemma}[theorem]{Lemma}
\newtheorem{proposition}[theorem]{Proposition}
\newtheorem{corollary}[theorem]{Corollary}
\newtheorem{example}[theorem]{Example}
\theoremstyle{definition}
\newtheorem{definition}[theorem]{Definition}
\newcommand{\R}{\mathbb R}
\newcommand{\Q}{\mathbb Q}
\title{Exact contracts on an affine line of outcome distributions}
\author{GPT 6 Astra Ultra}
\date{9 October 2026}
\begin{document}
\maketitle
\noindent\textbf{Problem D86-21. Status: unresolved.} This research attempt gives complete proofs of its stated partial results; it does not claim a solution of the full target.

\section{Recap and what is already known}
Action $a$ has rational cost $c_a\geq0$ and outcome distribution $q_a$;
one action has zero cost. The principal's rational outcome rewards $r_j$
are nonnegative. A contract $t\geq0$ pays $t_j$ at outcome $j$.
The agent maximizes $q_a\cdot t-c_a$, with ties favoring the principal.
The target asks for an optimal contract with a number of arithmetic operations
polynomial in the numbers $n,m$ of actions and outcomes, independent of numeric
bit lengths, and polynomial-size intermediate encodings.

The survey~\cite{survey}, Section 3.1, gives the minimum-payment LP approach,
ordinary polynomial-time solvability, and the at-most-$n-1$-nonzero-payments
observation. Those facts are not novel contributions of this attempt. We
rederive the support fact and develop an explicit strongly polynomial algorithm
when the distributions lie on an affine line. This restriction allows arbitrarily
many actions and outcomes but does not cover general technologies.

\section{Why sparse payments alone do not settle the target}
For a candidate action $a$, put $B_{bj}=q_{aj}-q_{bj}$ and
$d_b=c_a-c_b$ for $b\ne a$. Its minimum-payment problem is
\[
 \min q_a\cdot t\quad\text{subject to }Bt\geq d,\quad t\geq0.
\]
\begin{lemma}[The known support bound, with degeneracy allowed]
If action $a$ is weakly inducible, its minimum-payment problem has an optimal
solution with at most $\operatorname{rank}(B)\leq n-1$ positive coordinates.
\end{lemma}
\begin{proof}
The objective is nonnegative on the feasible set. A bounded-below linear
objective on a nonempty polyhedron attains its infimum; apply the standard
finite-dimensional LP attainment theorem. Among optima choose one, $t$, with
minimum support. If its supported columns of $B$ are dependent, choose a
nonzero vector $z$ supported there with $Bz=0$. For sufficiently small positive
and negative $\epsilon$, $t+\epsilon z$ remains nonnegative and feasible.
Optimality forces $q_a\cdot z=0$, since otherwise one sign improves the objective.
Choose a sign and increase its magnitude until the first positive coordinate
becomes zero. Such a finite endpoint exists in at least one direction because
$z\ne0$. This gives another optimum with strictly smaller support, a
contradiction. Thus the supported columns are independent.
\end{proof}

\begin{proposition}[Fixed number of actions]
For every fixed $n$, an optimal finite contract can be found by a strongly
polynomial algorithm whose arithmetic-operation count is
$m^{O(n)}\operatorname{poly}(n)$ and whose intermediate encodings are polynomial
in input length. This is not a uniform strongly polynomial bound when $n$ varies.
\end{proposition}
\begin{proof}
For each $a$, the feasible polyhedron lies in the nonnegative orthant and
therefore has no lines. A nonempty optimal face has a vertex. At a vertex
with $k$ positive coordinates, the active incentive inequalities restricted
to these coordinates must have rank $k$; otherwise a small two-sided feasible
perturbation contradicts being a vertex. In particular $k\leq n-1$.
Enumerate each support of size $k\leq n-1$ and each choice of $k$ incentive
rows; solve the resulting square system when nonsingular and check all original
inequalities and nonnegativity. Include $t=0$ separately. At least one optimal
vertex is encountered for every inducible action. Gaussian elimination on
matrices of size at most $n-1$ and all comparisons use a number of operations
bounded as stated. Cramer's-rule determinant bounds give polynomial input-bit
length for all candidates. Select the action/payment candidate maximizing
$q_a\cdot(r-t)$.

To justify the tie rule, every candidate action is a utility maximizer for its
contract. If another tied action would give the principal more, principal-favorable
tie breaking realizes at least the candidate payoff. Conversely every actual
contract/action outcome is feasible in the corresponding LP. Thus maximizing
candidate payoffs attains exactly the optimal realizable payoff.
\end{proof}

\section{An unrestricted-outcome but one-dimensional technology}
\begin{theorem}[Affine-line contract algorithm]
Suppose the distributions have a rational representation
\[
 q_a=q^0+\theta_a d,\qquad \sum_jd_j=0,
\]
where all $q_a$ are probability vectors. If $d\ne0$, an optimal contract can
be computed using $O(n^2+nm)$ arithmetic operations, with polynomial-size
intermediate rational encodings. At most one outcome receives a positive
payment. If all $q_a$ coincide, the zero contract is optimal.
The affine-line hypothesis itself can be tested, and such a representation
found if it exists, with $O(nm)$ arithmetic operations.
\end{theorem}
\begin{proof}
First take $d\ne0$ and set $x=d\cdot t$. The incentive inequalities for $a$
reduce exactly to
\[
 (\theta_a-\theta_b)x\geq c_a-c_b\quad\text{for every }b.
\]
Intersect these half-lines to obtain a closed interval $I_a$, possibly empty
or unbounded. A zero coefficient with a positive right side makes it empty.
This takes $O(n)$ operations for each $a$.
Since $d$ is nonzero and sums to zero, both $\{j:d_j>0\}$ and
$\{j:d_j<0\}$ are nonempty. Define
\[
 \alpha_a^+=\min_{d_j>0}\frac{q_{aj}}{d_j},\qquad
 \alpha_a^-=\min_{d_j<0}\frac{q_{aj}}{-d_j}.
\]
Both numbers are nonnegative. For $x\geq0$, every $t\geq0$ with $d\cdot t=x$
satisfies $q_a\cdot t\geq\alpha_a^+x$: for $d_j>0$ this holds termwise;
for $d_j\leq0$, $q_{aj}\geq0\geq\alpha_a^+d_j$.
Equality is achieved by selecting a minimizing outcome $j$ and setting
$t_j=x/d_j$, all other payments zero. For $x\leq0$ the analogous bound and
one-outcome construction give minimum payment $\alpha_a^-(-x)$.
In particular at $x=0$ the zero contract achieves minimum payment zero.
Thus the exact minimum-payment cost as a function of $x$ is
\[
 C_a(x)=\begin{cases}\alpha_a^+x,&x\geq0,\\
 \alpha_a^-(-x),&x\leq0.\end{cases}
\]
It is nondecreasing away from zero on either half-line. Consequently choose
$x_a=0$ if $0\in I_a$, the left endpoint if $I_a\subset(0,\infty)$, or the
right endpoint if $I_a\subset(-\infty,0)$. The endpoint is finite whenever
needed. This remains optimal if a slope vanishes. Realize $x_a$ by the
one-outcome payment described above, calculate
$q_a\cdot r-C_a(x_a)$, and maximize over nonempty $I_a$.
The preceding proposition's tie argument applies unchanged and proves global
optimality. There is at least one nonempty interval because zero payment
induces a minimum-cost action.

All intervals require $O(n^2)$ comparisons and divisions; all slope minima,
rewards, and output comparisons require $O(nm)$ more. Endpoints are ratios of
input costs and $\theta$ differences, and payments are products and ratios of
a bounded number of these rational quantities. Sums in rewards have at most
$m$ terms. Their encodings are polynomial in total input length.

For recognition, choose $q^0$ as one input row. If every row is equal, there
is no variation in expected payment or reward across actions. A minimum-cost
action is always chosen, and zero payment maximizes principal payoff.
Otherwise choose a nonzero row difference $d=q_b-q^0$, and one coordinate
$j_0$ with $d_{j_0}\ne0$. Set
$\theta_a=(q_{aj_0}-q^0_{j_0})/d_{j_0}$ and check every coordinate of
$q_a=q^0+\theta_a d$. This proves and tests the hypothesis using $O(nm)$
operations, and the same encoding bounds apply.
\end{proof}

\begin{example}[A checked instance with interior action]
Let the two outcome rewards be $(0,10)$ and let success probabilities and costs
of the three actions be $(0,0),(1/2,1),(1,3)$.
The minimum payments for the three actions are respectively $0,1,4$ in
expectation: inducing the middle action requires success-minus-failure payment
in $[2,4]$, and inducing the high action requires it at least $4$.
The corresponding maximum principal payoffs are $0,4,6$. A payment of $4$
on success and zero on failure is optimal; the high and middle actions give
the agent equal utility $1$, and principal-favorable tie breaking selects
the high action with principal payoff $6$.
\end{example}

\section{Status and first remaining gap}
The line hypothesis turns all incentives into constraints on one scalar and
makes minimum payments explicit. General rows have several independent
directions; neither the support bound nor fixed-$n$ enumeration yields an
operation count polynomial jointly in $n,m$. No reduction proving equivalence
to general strongly polynomial linear programming is asserted. The first gap
is a uniform method for those multidimensional incentive systems.

\paragraph{Finite implementation checks.}
As a finite algebraic check, the affine-line algorithm was compared with
complete vertex enumeration of every action's payment LP on 100 rational
four-action, three-outcome instances, using exact fraction arithmetic and
random seed $20261009$. Costs included a zero-cost action and rewards were
nonnegative integers. Every optimum and incentive-feasibility check agreed.
This test supports the implementation but does not replace the proof.

\begin{thebibliography}{9}
\bibitem{survey} P.~D\"utting, M.~Feldman, and I.~Talgam-Cohen,
\emph{Algorithmic Contract Theory: A Survey}, arXiv:2412.16384v1 (2024),
Section 3.1, Observations 3.3--3.4 and the intervening computational discussion;
Section 3.3 discusses simpler action and outcome domains.
\url{https://arxiv.org/html/2412.16384v1}.
\end{thebibliography}

\end{document}
